% S1.7: descriptive SHAP for the 2026 models (protocol 5e). Numbers from generated/.
We explain the two models the final run saved for this purpose, the boosted hazard and the sequence
neural hazard on blocks A--D, both fitted for the 2026 test year. The explained output is the
predicted probability that the swing reverses within five sessions. We used permutation SHAP on a
fixed random sample of \ShapRows{} rows of the 2026 test year, with \ShapBackground{} training rows
as the background. For the neural model each row's own 60-session price sequence is held at its
value, so its attributions are those of the static features given the sequence. The attributions
add up to each prediction (largest error \ShapGbmAdditivity{} for the boosted and
\ShapNnAdditivity{} for the neural model). This is a description of the fitted models, not a test.

Both models rely mostly on the swing state (Table~\ref{tab:shapblocks}): block A carries
\ShapGbmBlockA\% of the total attribution in the boosted model and \ShapNnBlockA\% in the neural
model, and the distance to the threshold alone carries \ShapGbmTopShare\% and \ShapNnTopShare\%
(Table~\ref{tab:shaptop}). The global block accounts for \ShapGbmBlockD\% and \ShapNnBlockD\%. The
picture matches the main results: the distance to the threshold dominates, and the other features
make small contributions, which the neural model spreads more evenly over the stock and local
market blocks.

\begin{table}[!htbp]
\centering
\caption{Share of the total mean absolute SHAP value by feature block, 2026 test year, two models
fitted for 2026, \ShapRows{} rows.}\label{tab:shapblocks}
\footnotesize
\gentable{tab_shap_blocks}
\end{table}

\begin{table}[!htbp]
\centering
\caption{The twelve features with the largest mean absolute SHAP value, in percentage points of
the five-session reversal probability ($\times 100$), with the feature's block in parentheses.
Dir.: sign of the correlation between the feature and its SHAP value (higher feature value, higher
or lower predicted probability).}\label{tab:shaptop}
\footnotesize
\gentable{tab_shap_top}
\end{table}
